\section{The unnormalized strong fixed-point relation}

Appendix~D of \cite[lines~2109--2116]{Cirac2016MPDO_arXiv} considers a
stronger local relation in which the two-site tensor is obtained by adjoining
a positive physical matrix to one copy of the tensor and applying a unitary
change of basis. The matrix is not assumed to have trace one.

\begin{definition}[Strong renormalization fixed-point relation]
    \label{def:mpdo_strong_rfp}
    \lean{MPOTensor.IsStrongRFP}
    \leanok
    \uses{def:mpo_tensor}
    An MPO tensor $M$ satisfies the strong renormalization fixed-point relation
    if it admits a matrix $P\geq0$ on the physical space and a unitary $U$ on
    the two-site physical space such that, for all two-site physical indices
    $i=(i_1,i_2)$ and $j=(j_1,j_2)$,
    \begin{align}
        M^{i_1j_1}M^{i_2j_2}
        &=\sum_{a,b}
          U_{i,a}P_{a_2b_2}\overline{U_{j,b}}\,M^{a_1b_1}.
        \label{eq:mpdo_strong_rfp_tensor}
    \end{align}
    Equivalently, the literal orientation of the source diagram is
    $M_2=U(M_1\otimes P)U^\dagger$, with $M$ on the left physical site and
    $P$ on the right. There is no virtual change of basis and no scalar
    proportionality. This is the relation labelled \emph{Strong-RFP} in
    \cite[Appendix~D, lines~2109--2113]{Cirac2016MPDO_arXiv}.
\end{definition}

\begin{theorem}[Physical closures of the strong fixed-point relation]
    \label{thm:mpdo_strong_rfp_phys_close}
    \lean{MPOTensor.strongRFP_tensor_relation_iff_physClose,
      MPOTensor.isStrongRFP_iff_physClose}
    \leanok
    \uses{def:mpdo_strong_rfp, def:phys_close}
    For fixed matrices $P$ and $U$, the tensor identity
    (\ref{eq:mpdo_strong_rfp_tensor}) holds if and only if every virtual
    matrix $X$ satisfies
    \begin{align}
        M_2(X)
        &=U\bigl(M_1(X)\otimes P\bigr)U^\dagger.
        \label{eq:mpdo_strong_rfp_closure}
    \end{align}
    Hence the strong renormalization fixed-point relation is equivalent to the
    existence of $P\geq0$ and a two-site unitary $U$ satisfying
    (\ref{eq:mpdo_strong_rfp_closure}) for every $X$.
\end{theorem}

\begin{proof}
    \leanok
    \uses{thm:trace_nondeg}
    Multiplying (\ref{eq:mpdo_strong_rfp_tensor}) on the right by $X$ and
    taking the virtual trace gives, entrywise,
    \begin{align}
        \tr\!\left(M^{i_1j_1}M^{i_2j_2}X\right)
        &=\sum_{a,b}U_{i,a}P_{a_2b_2}\overline{U_{j,b}}
          \tr\!\left(M^{a_1b_1}X\right),
        \notag
    \end{align}
    which is (\ref{eq:mpdo_strong_rfp_closure}). Conversely, subtract the
    right-hand side of (\ref{eq:mpdo_strong_rfp_tensor}). Equation
    (\ref{eq:mpdo_strong_rfp_closure}) says that its product with every $X$
    has trace zero. Nondegeneracy of the matrix trace pairing makes the
    difference zero.
\end{proof}

\begin{theorem}[Global operator identity for fixed strong fixed-point witnesses]
    \label{thm:mpdo_strong_rfp_global_operator_identity}
    \lean{MPOTensor.strongRFP_physCloseN_eq_localized_preparation,
      MPOTensor.strongRFP_physCloseN_eq_reindex_singleKrausMap_kronecker,
      MPOTensor.strongRFP_mpo_eq_localized_preparation,
      MPOTensor.strongRFP_mpo_eq_reindex_singleKrausMap_kronecker}
    \leanok
    \uses{thm:mpdo_strong_rfp_phys_close,
      def:finite_tuple_regroupings, def:mpdo_global_renormalization_maps,
      def:mpdo_state_preparation_map, def:mpdo_single_kraus_map}
    Fix witnesses $P$ and $U$ in
    (\ref{eq:mpdo_strong_rfp_closure}). For every $N\geq0$ and every virtual
    matrix $X$, the $(N+2)$-site closure is obtained from the $(N+1)$-site
    closure by adjoining $P$ to the first physical coordinate and conjugating
    the first two physical coordinates by $U$. More explicitly, after the
    canonical identifications
    \begin{align}
      (\mathbb C^d)^{\otimes(N+1)}
        &\cong \mathbb C^d\otimes(\mathbb C^d)^{\otimes N},
        \notag\\
      (\mathbb C^d)^{\otimes(N+2)}
        &\cong (\mathbb C^d\otimes\mathbb C^d)
          \otimes(\mathbb C^d)^{\otimes N},
        \notag
    \end{align}
    one has
    \begin{align}
      M_{N+2}(X)
        &=(U\otimes\Id)\bigl(M_{N+1}(X)\otimes P\bigr)
          (U^\dagger\otimes\Id).
        \label{eq:mpdo_strong_rfp_global_closure}
    \end{align}
    The same identity holds for the ordinary periodic operators
    $\rho^{(N)}(M)$.
\end{theorem}

\begin{proof}
    \leanok
    Localize the one-to-two map
    $Y\mapsto U(Y\otimes P)U^\dagger$ at the first site and use
    (\ref{eq:mpdo_strong_rfp_closure}) on each matrix slice. The canonical
    regroupings identify localized preparation with the Kronecker product by
    $P$, and localized unitary conjugation with conjugation by
    $U\otimes\Id$. Closing the virtual indices with the identity gives the
    periodic-operator statement.
\end{proof}

\begin{theorem}[Geometric periodic rank growth of a strong fixed point]
    \label{thm:mpdo_strong_rfp_periodic_rank_growth}
    \lean{MPOTensor.strongRFP_rank_mpo_add_two,
      MPOTensor.strongRFP_rank_mpo_add_one,
      MPOTensor.IsStrongRFP.exists_periodic_rank_factor}
    \leanok
    \uses{def:mpdo_strong_rfp,
      thm:mpdo_strong_rfp_global_operator_identity}
    Fix witnesses $P\geq0$ and $U$ for the strong fixed-point relation. Then
    the ordinary full periodic MPO matrix ranks satisfy
    \begin{align}
      \operatorname{rank}\rho^{(N+2)}(M)
        &=\operatorname{rank}\rho^{(N+1)}(M)\,
          \operatorname{rank}P
          \qquad(N\geq0),
        \label{eq:mpdo_strong_rfp_rank_recurrence}\\
      \operatorname{rank}\rho^{(N+1)}(M)
        &=\operatorname{rank}\rho^{(1)}(M)\,
          (\operatorname{rank}P)^N.
        \label{eq:mpdo_strong_rfp_rank_geometric}
    \end{align}
    Consequently every strong fixed point admits a single positive witness
    $P$ whose matrix rank is the factor in both formulas. No normalization of
    $P$ and no positivity or normalization assumption on the periodic
    operators is required.
\end{theorem}

\begin{proof}
    \leanok
    Unitary conjugation and relabelling of matrix coordinates preserve rank,
    while rank is multiplicative under Kronecker products. Applying these
    facts to (\ref{eq:mpdo_strong_rfp_global_closure}) gives the recurrence.
    Induction from chain length one gives the geometric formula.
\end{proof}
